\section{How the actual radial diffusion recovers hidden arithmetic axis data}
\label{sec:ns-hidden-mellin-diffusion}

The fixed-axis singularity of Section~\ref{sec:ns-axis-arithmetic-jet}
has a precise evolution law in the same arithmetic coordinates.
An even radial Taylor coefficient disappears from the Taylor series
of the arithmetic trace, but survives as a regular Mellin value at
a trivial zeta zero. Radial diffusion reads that value and sends it
into an odd endpoint derivative. We prove the full maps and their
inverses, rather than discarding the canceled factor or replacing the
Navier--Stokes system by a scalar heat equation.

\subsection{The original transform and its complete input class}

Fix the source parameter \(0<h<1/100\), without changing it, and put
\[
 \mathcal D_0=\{a\in C^\infty([0,\infty);\mathbb C):
                   \operatorname{supp}a\text{ compact},\ a(0)=0\}.
\]
The coordinate of \(a\) is the physical \(\sigma=r^2/2\) when the
input is a fluid field. The arithmetic argument is denoted \(x\).
Retain exactly
\[
 \phi_a(x)=a(x)-2^ha(2x),\qquad
 \psi_a(x)=\phi_a(x)-\tfrac12\phi_a(x/2),\qquad
 G_a(x)=\mathcal A_ha(x)=\sum_{n\ge1}\psi_a(nx).
\]
The mean of \(\psi_a\) is zero by substitution, and \(\psi_a(0)=0\).
Section~\ref{sec:ns-axis-arithmetic-jet} proves by differentiated
Euler--Maclaurin remainders that \(G_a\in\mathcal D_0\), including
its smooth extension at zero, and gives
\begin{equation}\label{ndm:jet}
 G_a^{(j)}(0)=W_h(-j)a^{(j)}(0)\quad(j\ge1),\qquad
 W_h(s)=\zeta(s)(1-2^{s-1})(1-2^{h-s}).
\end{equation}
This \(W_h\) is entire: its first dilation factor cancels the
simple pole of \(\zeta\) at one. In particular \(G_a^{(2m)}(0)=0\)
for every \(m\ge1\).

There has nevertheless been no loss of the full function. On the
actual range \(\mathcal G_h=\mathcal A_h\mathcal D_0\), the inverse is
\begin{equation}\label{ndm:inverse}
 \begin{split}
 \psi_a(x)&=\sum_{n\ge1}\mu(n)G_a(nx),\\
 \phi_a(x)&=\sum_{\ell\ge0}2^{-\ell}\psi_a(x/2^\ell),\\
 a(x)&=-\sum_{k\ge1}2^{-kh}\phi_a(x/2^k).
 \end{split}
\end{equation}
For the first identity both divisor sums are finite at a fixed
\(x>0\). The identity
\(\sum_{d\mid m}\mu(d)=1\) for \(m=1\), and zero otherwise,
follows by expanding \(\prod_{p\mid m}(1-1)\); it proves that inverse.
The partial second sum ends in
\(\phi_a(x)-2^{-\ell-1}\phi_a(x/2^{\ell+1})\), and the partial
third sum ends in \(a(x)-2^{-kh}a(x/2^k)\).
The remainders tend to zero by boundedness at zero and \(h>0\).
After \(j\) differentiations the extra factors are \(2^{-\ell j}\)
and \(2^{-kj}\), so the recovered series converge smoothly on compact
half-line intervals. These statements concern the recovered \(\psi_a\)
and \(\phi_a\); no uniform bound on the first inverse near \(x=0\)
has been assumed. The three formulas prove both inverse compositions
on the stated range, as in Section~\ref{sec:ns-angular-mellin}.

\subsection{Every input Taylor coefficient in residues and regular values}

Use ordinary Taylor coefficients \(a_n=a^{(n)}(0)/n!\), not derivatives
without their factorials. For \(a\in\mathcal D_0\), its Mellin
transform is initially holomorphic for \(\Re s>-1\).
For any positive \(b\) and integer \(N\ge1\), subtraction at zero gives
the exact continuation
\begin{equation}\label{ndm:continuation}
 \begin{split}
 \widehat a(s)={}&
 \int_0^b\left(a(\sigma)-\sum_{n=1}^Na_n\sigma^n\right)
                       \sigma^{s-1}\,d\sigma
 +\int_b^\infty a(\sigma)\sigma^{s-1}\,d\sigma\\
 &+\sum_{n=1}^N\frac{a_n b^{s+n}}{s+n},\qquad \Re s>-N-1 .
 \end{split}
\end{equation}
Taylor's integral remainder is \(O(\sigma^{N+1})\). Each complex
derivative in \(s\) adds an integrable power of \(\log\sigma\) on
compact subsets of this half-plane; the second integral is entire.
This proves holomorphy of the remainder integrals, not merely formal
continuation. Different \(N,b\) agree on the original half-plane
and then by the identity theorem. Thus \(\widehat a\) is meromorphic
on \(\mathbb C\), with at most simple poles at \(-1,-2,\ldots\), and
\(\operatorname{Res}_{s=-n}\widehat a=a_n\).
The same argument applies to the proved smooth compact \(G_a\).

The exact Mellin formula \eqref{nam:arithmetic}, initially proved by
absolute summation and then analytic continuation, therefore holds
meromorphically on the entire plane:
\begin{equation}\label{ndm:product}
 \mathfrak G_a(s):=\mathcal M G_a(s)=W_h(s)\widehat a(s).
\end{equation}
The necessary local factors of \(W_h\) are computable from the
classical functional equation
\[
 \zeta(s)=2(2\pi)^{s-1}\sin(\pi s/2)\Gamma(1-s)\zeta(1-s),
\]
with its original constants \cite[\S25.4, (25.4.2)]{apostolDLMF25}.
At \(s=-2m\), only the sine vanishes; differentiating it there proves
\begin{equation}\label{ndm:zeta-derivative}
 \zeta'(-2m)=
 \frac{(-1)^m(2m)!}{2(2\pi)^{2m}}\zeta(2m+1)\ne0.
\end{equation}
Indeed \(\cos(-m\pi)=(-1)^m\), \(\Gamma(1+2m)=(2m)!\),
and \(\zeta(2m+1)=\sum_{n\ge1}n^{-2m-1}>0\).
Both dilation factors at \(-2m\) are nonzero. At negative odd
integers the sine is itself nonzero, and the other factors of the
functional equation are finite and nonzero. Hence the two exact
recovery formulas are
\begin{equation}\label{ndm:all-coefficients}
 \boxed{\quad
 a_{2m-1}=
 \frac{\operatorname{Res}_{s=-(2m-1)}\mathfrak G_a(s)}
       {W_h(-(2m-1))},\qquad
 a_{2m}=\frac{\mathfrak G_a(-2m)}{W_h'(-2m)}
       \quad(m\ge1).\quad}
\end{equation}
For the second formula write, in the local coordinate \(w=s+2m\),
\[
 \widehat a(s)=a_{2m}/w+H(w),\qquad
 W_h(s)=w\{W_h'(-2m)+wV(w)\},
\]
where \(H,V\) are holomorphic. Their product is holomorphic and has
value \(W_h'(-2m)a_{2m}\) at \(w=0\).
This also proves the formula when \(a_{2m}=0\).
For odd indices the residue of a product with a holomorphic nonzero
factor gives the first formula, including the zero-residue case.
Every higher coefficient of the analytic unit has been retained in
the displayed local product; none changes its constant term.

For later use define
\begin{equation}\label{ndm:constants}
 \kappa_h=W_h(-1)=\frac{2^{h+1}-1}{16}>0,\qquad
 d_h=W_h'(-2)=
       \frac{7(2^{h+2}-1)\zeta(3)}{32\pi^2}>0.
\end{equation}
The first identity is the original \(\zeta(-1)=-1/12\) calculation.
The second follows from \eqref{ndm:zeta-derivative} with \(m=1\),
\(1-2^{-3}=7/8\), and the remaining negative dilation factor.
In particular
\[
 \frac{\kappa_h}{d_h}
 =\frac{2\pi^2(2^{h+1}-1)}
         {7(2^{h+2}-1)\zeta(3)} .
\]

Thus cancellation at a trivial zero removes a pole from
\(\mathfrak G_a\), but does not remove the input coefficient:
it becomes the regular value in \eqref{ndm:all-coefficients}.
The vanishing even Taylor coefficients of \(G_a\) and its generally
nonzero regular Mellin values are exact related coordinates, not
contradictory assertions. These recover all Taylor coefficients of
\(a\), not its whole smooth germ from that Taylor sequence alone;
smooth flat remainders are retained by the full inverse
\eqref{ndm:inverse}.

\subsection{The actual diffusion operator and a proved endpoint obstruction}

The radial part of the angular-momentum diffusion in
\eqref{nam:pde} is
\[
 L_{\rm rad}a(\sigma)=2\sigma a''(\sigma).
\]
It maps \(\mathcal D_0\) to itself. Its exact transported operator,
with domain and codomain both specified, is
\begin{equation}\label{ndm:transport}
 \mathscr L_h:\mathcal G_h\longrightarrow\mathcal G_h,\qquad
 \mathscr L_hG=\mathcal A_hL_{\rm rad}\mathcal A_h^{-1}G.
\end{equation}
The inverse in this definition is precisely the three convergent
series \eqref{ndm:inverse}. It is not division by \(W_h\) at its zeros.
For \(\Re s>0\), integration by parts twice, with the zero-end
terms vanishing since \(a=O(\sigma)\) and \(a'=O(1)\), gives
\[
 \mathcal M(L_{\rm rad}a)(s)=2s(s-1)\widehat a(s-1).
\]
Upper boundary terms vanish by compact support. It follows that
\begin{equation}\label{ndm:radial-mellin}
 \mathcal M(\mathscr L_hG_a)(s)
       =2s(s-1)W_h(s)\widehat a(s-1).
\end{equation}
This equality then holds meromorphically everywhere by
\eqref{ndm:continuation}. In particular
\begin{equation}\label{ndm:radial-residue}
 \boxed{\quad
 (\mathscr L_hG_a)'(0)
 =\operatorname{Res}_{s=-1}\mathcal M(\mathscr L_hG_a)(s)
 =4\kappa_h a_2
 =\frac{4\kappa_h}{d_h}\mathfrak G_a(-2).\quad}
\end{equation}
The direct derivative gives the same factor:
\((2\sigma a'')'(0)=2a''(0)=4a_2\), followed by
\eqref{ndm:jet} for its arithmetic image.

There is no operator on full Taylor sequences at \(x=0\) whose
output is the full Taylor sequence of \(\mathscr L_hG\) for every
\(G\in\mathcal G_h\). Here is an explicit witness, including its
smooth cutoff. Put
\[
 e(v)=
 \begin{cases}\exp(-1/v),&v>0,\\0,&v\le0,\end{cases}
 \qquad
 \chi(\sigma)=
 \frac{e(2-\sigma)}{e(2-\sigma)+e(\sigma-1)},\qquad \sigma\ge0.
\]
The denominator is strictly positive: both terms could vanish only
if \(\sigma\ge2\) and \(\sigma\le1\) simultaneously.
Every right derivative of \(e\) at zero is zero, because on \(v>0\)
each derivative is a polynomial in \(1/v\) times \(e^{-1/v}\), and
\(y^ke^{-y}\to0\) as \(y\to+\infty\). Thus \(\chi\) is smooth,
equals one on \([0,1]\), and is zero on \([2,\infty)\).
Take \(a(\sigma)=\sigma^2\chi(\sigma)\).
Its only nonzero endpoint derivative is \(a''(0)=2\).
Equation \eqref{ndm:jet} proves that \(G_a\) is flat at zero:
every one of its derivatives there is zero, including its value.
It is not the zero function, by injectivity of \(\mathcal A_h\).
But \eqref{ndm:radial-residue} gives
\[
 (\mathscr L_hG_a)'(0)=4\kappa_h\ne0,\qquad
 \mathscr L_h0=0.
\]
The two inputs \(G_a\) and zero have identical full Taylor sequences
and different output first derivatives. This proves the assertion
without presuming any finite-order description.
In particular \(\mathscr L_h\) cannot agree on its entire stated
domain with a finite-order differential operator having coefficients
smooth up to \(x=0\): repeated product rules show that every such
operator sends flat functions to flat functions.
The obstruction concerns endpoint jets, not injectivity of the
full arithmetic transform, whose exact inverse was proved above.

More generally,
\[
 (\mathscr L_hG_a)^{(j)}(0)
      =2jW_h(-j)a^{(j+1)}(0)\quad(j\ge1).
\]
This follows by applying \eqref{ndm:jet} to \(2\sigma a''\)
and differentiating the product. For odd \(j\) it accesses precisely
an even input derivative. The regular-value half of
\eqref{ndm:all-coefficients} supplies that derivative with its exact
factorial. No Taylor datum has been silently invented.

\subsection{The complete nonlinear residue equation for the actual source}

Use exactly the actual corrected-field functions of
Section~\ref{sec:ns-angular-mellin}, with their physical viscosity:
\[
 B=\langle r u_\theta\rangle,\quad
 P=\langle r^2u_ru_\theta\rangle,\quad
 Q=\langle r u_z u_\theta\rangle,\quad
 F=\langle r f_\theta\rangle,\qquad
 \langle\cdot\rangle=\frac1{2\pi}\int_0^{2\pi}(\cdot)\,d\theta .
\]
Their compact supports, all original cutoff-summed cross terms,
and their exact equation are already proved in \eqref{nam:pde}.
Write ordinary radial Taylor coefficients
\[
 B=b_1\sigma+b_2\sigma^2+O(\sigma^3),\qquad
 P=p_2\sigma^2+O(\sigma^3),\qquad
 Q=q_1\sigma+O(\sigma^2),\quad F=f_1\sigma+O(\sigma^2).
\]
These are coefficients of the complete actual angular averages;
no axisymmetry has been imposed outside the source's inner region.
The coefficient of \(\sigma\) in \eqref{nam:pde} is exactly
\begin{equation}\label{ndm:axis-balance}
 \partial_tb_1+2p_2+\partial_zq_1
       =\nu(4b_2+\partial_{zz}b_1)+f_1 .
\end{equation}
This follows equivalently by one \(\sigma\) derivative at zero;
the smoothness already proved justifies commutation with \(z,t\)
derivatives on every compact preterminal parameter set.

Let \(\mathfrak B=\mathcal M\mathcal A_hB\), and define
\(\mathfrak P,\mathfrak Q,\mathfrak F\) identically.
Define their residues
\[
 J_B=\operatorname{Res}_{s=-1}\mathfrak B,\qquad
 J_Q=\operatorname{Res}_{s=-1}\mathfrak Q,\qquad
 J_F=\operatorname{Res}_{s=-1}\mathfrak F .
\]
Equation \eqref{ndm:all-coefficients} proves, without deleting any
flux, that
\[
 J_B=\kappa_hb_1,\quad J_Q=\kappa_hq_1,\quad
 J_F=\kappa_hf_1,\quad
 \mathfrak B(-2)=d_hb_2,\quad
 \mathfrak P(-2)=d_hp_2.
\]
Multiplying \eqref{ndm:axis-balance} by \(\kappa_h\) therefore gives
the full arithmetic evolution identity
\begin{equation}\label{ndm:arithmetic-balance}
 \boxed{\quad
 \partial_tJ_B+\partial_zJ_Q-\nu\partial_{zz}J_B-J_F
 =\frac{\kappa_h}{d_h}
      \{4\nu\mathfrak B(-2)-2\mathfrak P(-2)\}.\quad}
\end{equation}
It holds on the entire preterminal domain where the actual source
fields are smooth. The right side retains the true quadratic flux
and radial diffusion, recovered by regular Mellin values. It is not
a closed scalar equation for \(J_B\).

The division-free equation \eqref{nam:arithmeticpde} yields exactly
the same identity, including at its canceled factor. In the local
coordinate \(w=s+1\), its left side is the product of
\[
 W_h(s-1)=d_hw+O(w^2)
 \quad\hbox{and}\quad
 \frac{\partial_tJ_B+\partial_zJ_Q-\nu\partial_{zz}J_B-J_F}{w}
            +O(1).
\]
Its limit is \(d_h\) times the numerator. Its right side tends to
\[
 -2\kappa_h\{\mathfrak P(-2)-2\nu\mathfrak B(-2)\}.
\]
Thus cancellation is a fully calculated zero-times-pole limit, not
the equation \(0=0\) obtained by substituting a zero factor before
retaining the principal part.

For the particular source field at \(z=0,t=1-\tau\),
\eqref{naj:residue} supplies the exact value
\[
 J_B(0,1-\tau)=\frac{2^{h+1}-1}{8C}\tau^{-1-h},
 \qquad
 \partial_tJ_B(0,1-\tau)
    =\frac{(1+h)(2^{h+1}-1)}{8C}\tau^{-2-h}.
\]
These can be substituted into \eqref{ndm:arithmetic-balance} without
discarding the transverse \(z\) derivative, the actual force, or
either regular Mellin value. The residue is the first radial
arithmetic derivative, as proved in Section~\ref{sec:ns-axis-arithmetic-jet}.
This identifies a specific part of its full dynamics: the even
radial curvature is invisible to the trace's Taylor series but
visible to its Mellin value at \(-2\), and enters the growing residue
through the original viscous term.
The auxiliary compact cutoff above proves a domain obstruction for
this exact operator; it is not a new Navier--Stokes or RH
counterexample. The actual source field and the trivial zeta zero
used in the evolution law have never been replaced.

\subsection{The full alternating residue--value hierarchy}

The same calculation retains every radial Taylor order, not just
the first growing derivative. For every integer \(n\ge1\) define
the nonzero constant and the corresponding Mellin coordinate
\[
 c_{h,n}=
 \begin{cases}W_h(-n),&n\text{ odd},\\W_h'(-n),&n\text{ even},\end{cases}
 \qquad
 E_n(\mathfrak G)=
 \begin{cases}
   \operatorname{Res}_{s=-n}\mathfrak G(s),&n\text{ odd},\\
   \mathfrak G(-n),&n\text{ even}.
 \end{cases}
\]
On the actual arithmetic image these values and residues exist by
\eqref{ndm:all-coefficients}, and
\(E_n(\mathfrak G_a)=c_{h,n}a_n\).
This is a diagonal bijection between the input Taylor coefficients
and their joint residue--value coordinates, with inverse
\(a_n=E_n/c_{h,n}\). It is not a claim that a smooth function is
determined by its Taylor coefficients.

For \(A=B,P,Q,F\), put
\(\widetilde A_n=E_n(\mathfrak A)\) and
\(A_n=\partial_\sigma^n A(0,z,t)/n!\).
Differentiating the exact equation \eqref{nam:pde} \(n\) times at
zero, or applying the product rule with these factorials, gives
\[
 \partial_tB_n+(n+1)P_{n+1}+\partial_zQ_n
 =\nu\{2n(n+1)B_{n+1}+\partial_{zz}B_n\}+F_n .
\]
For the radial diffusion coefficient, \(B_{n+1}\sigma^{n+1}\)
contributes \(2n(n+1)B_{n+1}\sigma^n\); every other monomial has
a different degree. Taylor's differentiated remainder justifies
this finite coefficient extraction without analyticity.
Multiplication by \(c_{h,n}\) and the proved coefficient inverse
therefore yield, at every order,
\begin{equation}\label{ndm:full-hierarchy}
 \begin{split}
 \partial_t\widetilde B_n+\partial_z\widetilde Q_n
       -\nu\partial_{zz}\widetilde B_n-\widetilde F_n
 = (n+1)\frac{c_{h,n}}{c_{h,n+1}}
             \{2\nu n\,\widetilde B_{n+1}
                             -\widetilde P_{n+1}\},\quad n\ge1 .
 \end{split}
\end{equation}
Every denominator has been proved nonzero from the functional
equation, with no RH assumption. This alternates regular values
at negative even integers with residues at negative odd integers.
The \(n=1\) equation is precisely \eqref{ndm:arithmetic-balance}.
All original nonlinear fluxes, forcing coefficients and viscosity
remain; the complete functions, including flat radial remainders,
are still carried by \eqref{ndm:product} and \eqref{ndm:inverse}.
